% JTDH 2026 submission — compile with pdfLaTeX + biber

\documentclass{article}
\usepackage{jtdh}
\usepackage[english]{babel}
\usepackage{url}
\usepackage[hang,flushmargin]{footmisc}
\usepackage{hyperref}
\hypersetup{colorlinks=true, linkcolor=black, urlcolor=gray, citecolor=black, anchorcolor=black}
\usepackage{booktabs}
\usepackage{amsfonts}
\usepackage{xcolor}
\usepackage{graphicx}
\usepackage{tabularx}
\usepackage[backend=biber, style=apa, uniquelist=false]{biblatex}
\addbibresource{jtdh_refs.bib}
\usepackage{enumitem}
\graphicspath{{img/}}
\aboverulesep=0ex
\belowrulesep=0ex

\title{Reading Without Having Read: How Reading Strategy, Interpretive Framework, and Textual Grounding Shape Computational Literary Discussion}

\author{
  Chris BREW\textsuperscript{1,\,2}\orcidlink{0000-0003-3946-4395}
}
\institutions{
  \textsuperscript{1}The Ohio State University, Columbus, Ohio, USA\par
  \textsuperscript{2}LexisNexis Legal \& Professional
}

\begin{document}
\maketitle

\begin{abstract}
A reader who surveys a novel broadly and a reader who studies it
narrowly will assemble different evidence for discussion.  Yet when we
implement these two reading strategies computationally and generate
multi-voice literary discussions about 15 novels (60,774 passages, 180
experimental conditions), the resulting conversations are
near-identical in character: expert airtime is stable to $\pm$2\%
regardless of strategy.  The two algorithms agree on fewer than 6\% of
their selections, yet the interpretive framework---not the reading
strategy---determines what gets said.  This convergence provides
computational evidence for \textcite{fish1980text}'s thesis that
meaning is reader-constructed.  Removing textual grounding produces not
random hallucination but \emph{measurable pastiche}: 82\% of invented
quotes draw 90--100\% of their vocabulary from the novel's word stock,
with fidelity ranging from 57\% (\emph{Middlemarch}) to 0\%
(\emph{Hester}).  This extends \textcite{bayard2007comment}'s
philosophical argument about non-reading into an empirically testable
framework---and connects it to recent work on LLM memorisation of
training data \parencite{carlini2021extracting, carlini2023quantifying}.
\end{abstract}
\keywordseng{literary interpretation, reader-response theory, LLM memorisation, passage grounding, digital humanities}


%% ============================================================
\section{Two Ways of Reading, One Discussion}
\label{sec:intro}

Consider two scholars preparing to lead a seminar on
\emph{Middlemarch}.  The first reads broadly, marking passages across
thirty chapters that develop character and advance the plot---a
\emph{synoptic} reader.  The second reads selectively, concentrating
on passages dense with social critique and thematic argument---an
\emph{analytical} reader.  Both are legitimate strategies.  They will
arrive at the seminar with different evidence.  The question is whether
they will arrive at different discussions.

We can now answer this question computationally.  We implement the two
strategies as passage-selection algorithms applied to 15 Victorian and
early-modern novels, then generate structured multi-voice literary
discussions from each selection.  The synoptic method (min-cost network
flow; \nptextcite{ahuja1993network}) selects 59 passages per run from 28
chapters, favouring character development and plot advancement.  The
analytical method (contextual embeddings with LLM-curated similarity)
selects 32 passages from 20 chapters, favouring social critique and
thematic depth.  The two agree on fewer than 6\% of their selections
(Jaccard 0.057).

Yet the resulting discussions converge.  Each expert's share of
airtime varies by less than 2 percentage points regardless of which
algorithm did the selecting.  The interpretive framework---embodied in
expert personas with distinct critical orientations---dominates.  The
reading strategy does not.

This \emph{framework dominance} is the central finding of this paper.
It provides large-scale computational evidence for
\textcite{fish1980text}'s thesis that meaning is constituted by
interpretive communities, not extracted from texts.  We build on it in
three directions:

\begin{enumerate}
\item We show that convergence holds across 15 novels, with
interpretive frameworks maintaining stable character regardless of
textual affordance (Section~\ref{sec:convergence}).

\item We show that removing grounding produces \emph{measurable
pastiche}---extending \textcite{bayard2007comment}'s philosophical
argument about non-reading into an empirically testable
framework, and connecting it to research on LLM memorisation
\parencite{carlini2021extracting} (Section~\ref{sec:grounding}).

\item We show that conversational scaffolding transforms monologue
into dialogue uniformly across all novels, independently of content
(Section~\ref{sec:scaffolding}).
\end{enumerate}


%% ============================================================
\section{From Reader-Response Theory to Computational Experiment}
\label{sec:theory}

\subsection{Interpretive communities}

\textcite{fish1980text} argued that texts do not contain determinate
meanings; meaning is constituted by the interpretive strategies readers
bring.  Our framework dominance finding extends this: not only different
interpretive strategies but different \emph{reading
strategies}---synoptic vs.\ analytical---produce convergent
discussions, so long as the interpretive framework is held constant.

\subsection{Textual affordance}

\textcite{iser1978act} proposed that texts contain gaps that readers
fill according to their resources.  We operationalise affordance
through seven provision dimensions applied to every passage.  These
dimensions vary across novels, confirming that different texts afford
different interpretive possibilities---even as the framework dominance finding
shows that frameworks override affordances.

\subsection{Non-reading as a spectrum}

\textcite{bayard2007comment} argued provocatively that ``non-reading''
is not a binary state but a spectrum: books skimmed, books heard about,
books forgotten.  He suggested that cultural positioning matters more
than textual knowledge.  Our no-passages condition---in which the LLM
generates literary discussion from prior knowledge alone---provides the
first computational test of this claim.  The results confirm Bayard's
intuition, but with a crucial empirical twist: we can \emph{measure}
the graduated degradation from reading to non-reading, and show that it
tracks the model's training-data exposure to each novel.

\subsection{LLM memorisation and literary knowledge}

\textcite{carlini2021extracting} demonstrated that language models
memorise and can reproduce verbatim excerpts from training data.
\textcite{carlini2023quantifying} showed that memorisation grows
log-linearly with model capacity, data duplication, and prompt length.
\textcite{duarte2024decop} developed methods to detect which books are
in a model's training data.  Our pastiche finding connects this
technical literature to literary-critical concerns: the LLM's
differential ability to quote different novels is a measurable trace of
its training-data composition, visible as a gradient from faithful
quotation of canonical works to pure pastiche of obscure ones.


%% ============================================================
\section{System and Experimental Design}
\label{sec:system}

We generate structured literary discussion podcasts through a pipeline
that decomposes interpretation into manipulable phases.  The system is
the experimental instrument; a companion technical paper describes
implementation details.

\paragraph{Corpus.}  Fifteen novels (Table~\ref{tab:novels}), parsed
into passages of 50--80 words.  Total: 60,774 passages across seven
authors (Dickens, Eliot, Gaskell, Forster, Collins, Gissing, Oliphant)
and seven decades (1850--1924).

\begin{table}[t]
\caption{Corpus: 15 novels, 60,774 passages, ordered by ungrounded
verification rate (Section~\ref{sec:grounding}).}
\label{tab:novels}
\begin{tabularx}{\textwidth}{X|l|r|r|r}
\toprule
\textit{Novel} & \textit{Author} & \textit{Year} & \textit{Passages} & \textit{Nop\%} \\
\midrule
Middlemarch          & Eliot    & 1871 & 5,519 & 57 \\
Bleak House          & Dickens  & 1853 & 6,916 & 53 \\
David Copperfield    & Dickens  & 1850 & 5,621 & 44 \\
Hard Times           & Dickens  & 1854 & 1,773 & 39 \\
A Passage to India   & Forster  & 1924 & 2,556 & 39 \\
Mill on the Floss    & Eliot    & 1860 & 4,654 & 37 \\
Cranford             & Gaskell  & 1853 & 1,226 & 33 \\
Daniel Deronda       & Eliot    & 1876 & 5,170 & 30 \\
Our Mutual Friend    & Dickens  & 1865 & 5,775 & 29 \\
New Grub Street      & Gissing  & 1891 & 3,872 & 20 \\
The Odd Women        & Gissing  & 1893 & 3,395 &  6 \\
Miss Marjoribanks    & Oliphant & 1866 & 3,019 &  3 \\
North and South      & Gaskell  & 1855 & 2,942 &  2 \\
Hester               & Oliphant & 1883 & 2,279 &  0 \\
No Name              & Collins  & 1862 & 6,057 &  0 \\
\bottomrule
\end{tabularx}
\end{table}

\paragraph{Two reading strategies.}  The synoptic strategy (min-cost
flow) selects 59 passages from 28 chapters, favouring
\texttt{character\_development} (+8.3~pp) and
\texttt{plot\_advancement} (+10.2~pp).  The analytical strategy
(embeddings) selects 32 from 20 chapters, favouring
\texttt{social\_critique} (+13.4~pp) and \texttt{thematic\_depth}
(+6.9~pp).

\paragraph{Experimental matrix.}  180 conditions: 15 novels $\times$
2 expert panels $\times$ 3 pipelines (synoptic, analytical,
no-passages) $\times$ 2 host-preparation settings (with/without).


%% ============================================================
\section{The Convergence Paradox}
\label{sec:convergence}

\begin{table}[t]
\caption{Synoptic vs.\ analytical reading across 28 matched pairs.}
\label{tab:convergence}
\begin{tabularx}{\textwidth}{X|r|r}
\toprule
\textit{Measure} & \textit{Synoptic} & \textit{Analytical} \\
\midrule
Passages per run           & 59    & 32    \\
Chapters per run           & 28    & 20    \\
Mean interest score        & 4.44  & 4.01  \\
\texttt{character\_development} strong & 83.5\% & 75.2\% \\
\texttt{plot\_advancement} strong      & 42.6\% & 32.5\% \\
\texttt{social\_critique} strong       & 50.6\% & 64.0\% \\
\texttt{thematic\_depth} strong        & 84.0\% & 90.9\% \\
\addlinespace
Passage Jaccard          & \multicolumn{2}{c}{0.057 (range 0.00--0.12)} \\
Expert airtime SD        & \multicolumn{2}{c}{$<$2~pp across 15 novels} \\
\bottomrule
\end{tabularx}
\end{table}

The two strategies select nearly disjoint passages with systematically
different provision profiles (Table~\ref{tab:convergence}).  The
synoptic reader covers the novel broadly; the analytical reader
concentrates on its arguments.

Yet expert airtime---each expert's share of total words, excluding the
host---is stable across all 15 novels and both strategies: Eleanor
Hartley (literary critic) 32--37\%, James Blackstone (social historian)
30--37\%, Caroline Woodcourt (close reader) 29--34\%.

This stability was not our prediction.  We expected expert prominence
to \emph{adapt} to textual affordance.  Instead, the frameworks are
\emph{robust}---stable heuristics for reading, not chameleons.  In
\textcite{fish1980text}'s terms: the reader makes the reading, and the
reading strategy does not.


%% ============================================================
\section{Reading Without Having Read}
\label{sec:grounding}

\subsection{The grounding gap}

With passages, the system quotes accurately: 98.0\% (synoptic), 97.2\%
(analytical).  Without passages, the mean drops to 48.9\%, ranging from
57\% (\emph{Middlemarch}) to 0\% (\emph{Hester}, \emph{No Name}).
This grounding gap ranges from $\sim$30~pp for canonical novels to
$\sim$100~pp for obscure ones.

\subsection{Bayard's spectrum, measured}

\textcite{bayard2007comment} argued that we always talk about
``approximate recollections of books, arranged as a function of our
current circumstances.''  The no-passages condition makes this literal:
the LLM's ``recollection'' of each novel is whatever it absorbed during
training.  Our verification rates measure this recollection on a
per-novel basis, producing a \emph{knowledge gradient} that tracks---but
does not perfectly follow---novel prominence.

\emph{Middlemarch} (57\% verified) and \emph{Bleak House} (53\%) are
canonical texts whose passages the model can often reproduce from
memory.  \emph{Hester} \parencite{oliphant1883hester} and \emph{No
Name} (both 0\%) are texts the model has barely encountered.  The
gradient between these extremes is itself a measurement of
literary-cultural prominence as reflected in training data.

\subsection{Pastiche, not hallucination}

We performed a vocabulary autopsy on 726 invented quotes (0\% match
with any passage):

\begin{itemize}
\item 82\% draw 90--100\% of their words from the novel's vocabulary;
\item mean overlap is 95.1\%; minimum is 67\%.
\end{itemize}

The LLM writes \emph{in the style of} the novel, not \emph{from} it.
It recombines real vocabulary into plausible-but-nonexistent sequences.
In literary terms, this is \emph{pastiche}.

This connects to \textcite{carlini2021extracting}'s finding that LLMs
memorise training data differentially.
\textcite{carlini2023quantifying} showed memorisation grows with data
duplication---and canonical novels are duplicated far more than obscure
ones in web-scraped training corpora.  Our pastiche gradient is the
literary-critical face of this technical phenomenon: the model's
ability to quote \emph{Middlemarch} reflects the frequency with which
\emph{Middlemarch} appeared in its training data.

\textcite{bender2021dangers} warned that language models are
``stochastic parrots'' producing text without understanding.  Our
pastiche results give this metaphor empirical texture: the parrot's
fidelity depends on how many times it has heard each text.


%% ============================================================
\section{Conversational Scaffolding}
\label{sec:scaffolding}

Host preparation (pre-interviews + question planning) raises questions
per segment from $\sim$1 to $\sim$11, \textbf{uniformly across all 15
novels}.  The ratio ranges from 5.3$\times$ (\emph{Our Mutual Friend})
to 14.1$\times$ (\emph{Cranford}).  Just as reading strategy does not
determine analytical character, literary content does not determine
conversational dynamics.  Both are shaped by framework---interpretive
in one case, architectural in the other.


%% ============================================================
\section{Weaknesses and Honest Limitations}
\label{sec:weaknesses}

We owe the reader a frank accounting of what this framework does
\emph{not} demonstrate, alongside what it does.

\subsection{LLM prompts are not a theory of interpretation}

Our ``interpretive frameworks'' are natural-language persona
descriptions fed to an LLM.  They produce outputs that
\emph{resemble} literary criticism performed from specific
perspectives.  But resemblance is not identity.  As
\textcite{shanahan2024talking} argues, we should resist the temptation
to map LLM behaviour onto human cognitive categories without
justification.  When we say the system ``performs a Marxist reading,''
we mean that its output contains vocabulary, framing, and emphasis
characteristic of Marxist criticism---not that it executes the
intellectual operations a Marxist critic would.  The convergence
paradox shows that persona prompts produce stable, distinctive
outputs.  It does \emph{not} show that these outputs are produced by
the same mechanisms as human literary interpretation.

To put it plainly: our personas are stereotypes, not theories.  They
are useful experimental instruments---they produce controlled,
reproducible variation in a generated text---but they should not be
mistaken for cognitive models of how human readers interpret
literature.  \textcite{fish1980text}'s interpretive communities are
social and cognitive phenomena; our expert personas are prompt
engineering.  The convergence between them---that both produce
framework-dominant, text-secondary interpretive behaviour---is
suggestive but not probative.

\subsection{Automated metrics only}

We report verification rates, airtime distributions, and structural
statistics, but no human evaluation of discussion quality, insight, or
literary value.  A discussion that scores well on our metrics could
still be intellectually shallow.  Human evaluation is the most
important missing piece.

\subsection{English canon, 1850--1924}

All 15 novels are English, from a narrow period.  We cannot claim
generalisability to other languages, periods, or genres.  The
provision dimensions reflect Victorian Realist priorities
(\texttt{atmosphere\_setting} fails to transfer), and a schema
developed from different traditions would capture different affordances.

\subsection{No causal claims}

Novel prominence correlates with verification rate, but we do not
establish the causal mechanism.  The correlation could reflect
training-data duplication \parencite{carlini2023quantifying}, prose
distinctiveness, or quotability---we cannot distinguish these.


%% ============================================================
\section{Related Work}
\label{sec:related}

\textcite{moretti2013distant} and \textcite{underwood2019distant}
established computational methods in literary studies through topic
modelling, network analysis, and stylometry.  Our work differs in
generating \emph{discourse about} literature and in operationalising
reader-response concepts.

Google's NotebookLM \parencite{notebooklm2024} generates audio
discussions from documents.  Our contribution is the explicit
decomposition into manipulable components, enabling the controlled
comparisons that yield the framework dominance finding.

\textcite{lewis2020rag} and \textcite{gao2024rag} establish
retrieval-augmented generation as the standard approach to grounding.
Our finding that retrieval \emph{method} matters far less than the
\emph{fact} of retrieval extends this into a domain where source--output
relationships are transparent.


%% ============================================================
\section{Conclusion}
\label{sec:conclusion}

A synoptic reader and an analytical reader, given the same novel and
the same interpretive framework, produce discussions whose character is
determined by the framework, not the reading strategy.  This
convergence---across 15 novels and 180 conditions---provides
computational evidence, though not proof, for
\textcite{fish1980text}'s thesis that meaning is reader-constructed.

When grounding is removed, the system produces graduated pastiche:
plausible quotations built from the novel's vocabulary, more faithful
for canonical novels and degrading to invention for obscure ones.  This
pastiche is the literary-critical face of differential LLM memorisation
\parencite{carlini2021extracting}, and gives empirical texture to
\textcite{bayard2007comment}'s argument that we are always talking
about ``approximate recollections'' of books.

We close with a caution.  These findings are suggestive, not
dispositive.  LLM personas are not theories of mind.  Automated
metrics are not literary judgements.  But the framework dominance finding, the
pastiche gradient, and the uniformity of scaffolding effects are
empirical regularities that any account of computational literary
discussion must accommodate.  The reader makes the reading.  The
reading strategy does not.


\acknowledgmenteng{
The author thanks the anonymous reviewers for their comments.
}

\urlstyle{same}
\printbibliography

\newpage
% Appendix: standalone or \input into main paper
% Requires: booktabs, tabularx, amsmath, hyperref, biblatex with jtdh_refs.bib

\section*{Appendix: Why Can the LLM Quote Middlemarch but Not Hester?}
\addcontentsline{toc}{section}{Appendix: Analytical Reinforcement}
\label{app:reinforcement}

\noindent
When asked to discuss a novel without source passages, the LLM produces
verifiable quotations at rates ranging from 57\%
(\emph{Middlemarch}) to 0\% (\emph{Hester}).  All 15 novels are
available from Project Gutenberg\footnote{%
  \url{https://www.gutenberg.org}.  Gutenberg texts are a standard
  component of web-scraped training corpora such as C4
  and The Pile.  We assume all 15 novels are present in the training
  data of the model used (Claude, Anthropic), though we cannot verify
  this for any proprietary model.}
and almost certainly appear in the model's training data.  The gradient
must therefore reflect something other than raw textual presence.  This
appendix investigates what that something is, subjects the
investigation to systematic critique, and states what survives.

\subsection*{A.1\quad Hypothesis: Analytical Reinforcement}

\textcite{carlini2023quantifying} showed that LLM memorisation grows
log-linearly with data duplication: sequences that appear more
frequently in training data are reproduced more faithfully.  For
novels, the Gutenberg source text appears once.  But for canonical
novels, specific passages are \emph{duplicated} across thousands of
secondary sources: study guides that excerpt key quotations, academic
papers that discuss them, Wikipedia articles that summarise plots with
embedded excerpts, blog posts, teaching materials, and social media.

A famous passage---the fog opening of \emph{Bleak House}, the pier
glass passage of \emph{Middlemarch}---may appear not only in the
Gutenberg text but across hundreds of such sources.  A passage from
Chapter~12 of \emph{Hester} appears on Gutenberg and virtually
nowhere else.  Our hypothesis: the LLM's per-novel quote fidelity
tracks the volume of this \textbf{analytical reinforcement}, not the
presence of the source text.

\subsection*{A.2\quad Data Sources and Composite Score}

We constructed a composite attention score from three proxies, each
capturing a different facet of a novel's analytical footprint:

\begin{enumerate}[leftmargin=2em]

\item \textbf{Goodreads ratings count.}  Source: Goodreads.com, accessed
March 2026.  This measures present-day readership.
Ratings counts range from 548 (\emph{Hester}) to 261,000
(\emph{David Copperfield}).  Log-normalised to $[0, 100]$ to reduce
the influence of extreme popularity:
$\mathrm{score} = 100 \times (\log r - \log r_{\min}) / (\log r_{\max} - \log r_{\min})$
where $r$ is the ratings count.

\item \textbf{Wikipedia article word count.}  Source: English Wikipedia,
accessed March 2026.  Word count of the novel's Wikipedia article,
measured by selecting all text content.  This measures encyclopaedic
and analytical attention: longer articles contain more plot summary,
critical discussion, and quoted excerpts.  Ranges from $\sim$850
(\emph{Miss Marjoribanks}) to $\sim$19,000 (\emph{David
Copperfield}).  Linearly normalised to $[0, 100]$.

\item \textbf{Study guide presence.}  Binary indicator: does the novel
have a SparkNotes, LitCharts, CliffsNotes, or Shmoop study guide?
Source: manual check of each platform, March 2026.  Study guides are
particularly relevant because they systematically excerpt and re-quote
key passages, producing exactly the kind of passage-level duplication
that \textcite{carlini2023quantifying}'s framework predicts will
strengthen memorisation.  Coded as 33.3 (present) or 0 (absent).
\end{enumerate}

The composite score is the unweighted mean of the three normalised
components, producing a scale of 0--100.  We acknowledge this scoring
is ad hoc: the choice of proxies, normalisation methods, and equal
weighting is not theoretically derived.  Different reasonable choices
would produce different numbers.  We report it as one defensible
operationalisation, not as a validated instrument.

\subsection*{A.3\quad Results}

Table~\ref{tab:attention} presents the full data.

\begin{table}[ht]
\centering
\small
\caption{Critical attention metrics and ungrounded (no-passages) quote
verification rates.  Data sources: Goodreads.com, English Wikipedia,
SparkNotes/LitCharts (all accessed March 2026).  Nop\% is the
percentage of LLM-generated quotes verified against the full enriched
passage corpus using 5-tier matching ($\geq$60\% = verified).}
\label{tab:attention}
\begin{tabularx}{\textwidth}{X|l|r|r|r|c|r|r}
\toprule
\textit{Novel} & \textit{Author} & \textit{Year} &
  \textit{GR ratings} & \textit{Wiki words} & \textit{SG} &
  \textit{Attn.} & \textit{Nop\%} \\
\midrule
Middlemarch        & Eliot    & 1871 & 181,000 & 8,750  & Y &  79.2 & 57 \\
Bleak House        & Dickens  & 1853 &  98,000 & 8,750  & Y &  75.9 & 53 \\
David Copperfield  & Dickens  & 1850 & 261,000 & 19,000 & Y & 100.0 & 44 \\
Hard Times         & Dickens  & 1854 &  75,000 & 6,750  & Y &  70.8 & 39 \\
Passage to India   & Forster  & 1924 &  73,000 & 4,350  & Y &  66.2 & 39 \\
Mill on the Floss  & Eliot    & 1860 &  59,000 & 3,000  & Y &  62.6 & 37 \\
Cranford           & Gaskell  & 1853 &  48,000 & 2,900  & N &  27.9 & 33 \\
Daniel Deronda     & Eliot    & 1876 &  27,000 & 3,500  & Y &  59.3 & 30 \\
Our Mutual Friend  & Dickens  & 1865 &  32,000 & 8,750  & Y &  69.8 & 29 \\
New Grub Street    & Gissing  & 1891 &   6,300 & 3,200  & N &  17.5 & 20 \\
The Odd Women      & Gissing  & 1893 &   5,500 & 1,200  & N &  13.1 &  6 \\
Miss Marjoribanks  & Oliphant & 1866 &   2,600 &   850  & N &   8.4 &  3 \\
North and South    & Gaskell  & 1855 & 182,000 & 8,750  & Y &  79.2 &  2 \\
Hester             & Oliphant & 1883 &     548 & 1,300  & N &   0.8 &  0 \\
No Name            & Collins  & 1862 &   7,900 & 3,200  & N &  18.7 &  0 \\
\bottomrule
\end{tabularx}
\end{table}

\paragraph{Overall correlation.}  Pearson $r = 0.72$, Spearman
$\rho = 0.70$ ($n = 15$).  Excluding \emph{North and South}
(discussed below): Pearson $r = 0.89$, Spearman $\rho = 0.90$
($n = 14$).

\subsection*{A.4\quad The North and South Outlier}

\emph{North and South} has an attention score of 79.2---comparable
to \emph{Middlemarch} and \emph{Bleak House}---but a nop verification
rate of only 2\%.  Its Goodreads count (182,000) is almost certainly
inflated by the 2004 BBC television adaptation, which generated a
large modern fandom that discusses the adaptation's casting, romance,
and divergences from the novel---not the novel's prose passage by
passage.

This suggests that the composite attention score, which treats
Goodreads popularity as equivalent to Wikipedia coverage and study
guide presence, conflates \emph{readership} with \emph{analytical
engagement}.  A novel can be widely read (or widely watched) without
generating the passage-level discussion that reinforces specific
textual fragments in training data.  \emph{Cranford} (also adapted by
the BBC, 2007) shows a milder version of this effect but remains
closer to the trend line.

We report the outlier-excluded correlation ($r = 0.89$) alongside the
full-sample correlation ($r = 0.72$), but we do not rely on the
exclusion.  The more important evidence is the within-author analysis
below.

\subsection*{A.5\quad Within-Author Gradients}

The most damaging objection to the analytical reinforcement hypothesis
is that the correlation might reflect \emph{author effects}: perhaps
Dickens's prose is inherently more memorable (whether by LLMs, by people or by both) than Oliphant's,
regardless of secondary discussion.
We address this by examining
within-author gradients.

\begin{table}[ht]
\centering
\small
\caption{Within-author nop verification gradients.  We assume that novels by the
same author share significant stylistic properties; author effects are mitigated.}
\label{tab:within_author}
\begin{tabularx}{\textwidth}{l|X|r|r}
\toprule
\textit{Author} & \textit{Novel} & \textit{Attn.} & \textit{Nop\%} \\
\midrule
Dickens  & Bleak House        & 75.9 & 53 \\
         & David Copperfield  & 100.0 & 44 \\
         & Hard Times         & 70.8 & 39 \\
         & Our Mutual Friend  & 69.8 & 29 \\
\addlinespace
Eliot    & Middlemarch        & 79.2 & 57 \\
         & Mill on the Floss  & 62.6 & 37 \\
         & Daniel Deronda     & 59.3 & 30 \\
\addlinespace
Gaskell  & Cranford           & 27.9 & 33 \\
         & North and South    & 79.2 &  2 \\
\addlinespace
Gissing  & New Grub Street    & 17.5 & 20 \\
         & The Odd Women      & 13.1 &  6 \\
\addlinespace
Oliphant & Miss Marjoribanks  &  8.4 &  3 \\
         & Hester             &  0.8 &  0 \\
\bottomrule
\end{tabularx}
\end{table}

\textbf{Dickens} (4 novels): Bleak House (53\%) $>$ David Copperfield
(44\%) $>$ Hard Times (39\%) $>$ Our Mutual Friend (29\%).  This
within-Dickens ranking tracks canonical status: Bleak House is the
most taught and discussed; Our Mutual Friend, though artistically
admired, is the least widely assigned.  Same author, plausibly similar prose style,
different critical prominence, different quote fidelity.

\textbf{Eliot} (3 novels): Middlemarch (57\%) $>$ Mill on the Floss
(37\%) $>$ Daniel Deronda (30\%). It is plausible that Middlemarch is far more widely
taught and discussed than Daniel Deronda.

\textbf{Gaskell} (2 novels): Cranford (33\%) $\gg$ North and South
(2\%).  This is the most striking within-author contrast and the one
we understand least well.  North and South has more study guide
coverage (GradeSaver, LitCharts, Course Hero, SuperSummary) and a
higher attention score (79.2 vs.\ 27.9) than Cranford, yet Cranford
dramatically outperforms.  North and South's modern Goodreads
popularity is inflated by the 2004 BBC adaptation, which may generate
discussion of the \emph{adaptation} without reinforcing the
\emph{novel's} specific passages.  But we cannot rule out other
explanations: Cranford's shorter length may make it more quotable per
passage, or its comic tone may produce more distinctive (and therefore
more memorisable) phrasing.  This pair is genuinely anomalous and
resists the simple analytical-reinforcement explanation.

\textbf{Gissing} (2 novels): New Grub Street (20\%) $>$ The Odd Women
(6\%).  New Grub Street is the better-known novel.

\textbf{Oliphant} (2 novels): Miss Marjoribanks (3\%) $\approx$
Hester (0\%).  Both near zero, consistent with Oliphant's general
absence from modern critical discussion.

\medskip\noindent
Within-author gradients exist for all five authors with multiple novels
and, with the exception of the Gaskell pair (discussed in A.4), they
track critical-cultural prominence.  \textbf{This is the strongest
evidence for the analytical reinforcement hypothesis}, because it
controls for the most plausible confound: author prose style.

\subsection*{A.6\quad Objections and Qualifications}

We subjected the thesis to systematic critique.  The principal
objections, and our responses:

\begin{description}[leftmargin=1.5em, labelindent=0em]

\item[Training data opacity.]  We do not know the exact composition of
the model's training data.  The attention score is a proxy for likely
duplication in web-scraped corpora, not a measurement of actual
duplication.  The hypothesis is consistent with
\textcite{carlini2023quantifying}'s findings but cannot be verified
for proprietary models.

\item[Ad hoc scoring.]  The composite score is one defensible
operationalisation among several.  We tested robustness informally:
dropping any single proxy changes $r$ by less than 0.1.  A more
principled instrument would require validated cultural-prominence
metrics that do not currently exist.

\item[Small sample.]  $n = 15$ supports trend identification but not
confident statistical inference.  The correlation ($r = 0.72$) has a
wide confidence interval.  We rely on the within-author evidence
(Table~\ref{tab:within_author}) as much as on the aggregate
correlation.

\item[Readership vs.\ analytical engagement.]  The North and South
outlier demonstrates that popular readership does not entail the kind
of passage-level analytical engagement that would reinforce specific
fragments in training data.  A more refined metric would distinguish
passage-quoting activity (study guides, academic close readings) from
generic discussion (reviews, adaptation commentary).  We have not
measured this distinction directly.

\item[The Bayard parallel is a metaphor.]
\textcite{bayard2007comment}'s argument concerns human cultural
navigation; our findings concern statistical weight distributions in a
language model.  The parallel---both humans and LLMs navigate literary
culture through graduated approximation---is evocative but the
mechanisms differ entirely.  We flag it as an analogy, not an identity.
As \textcite{shanahan2024talking} argues, care is needed in mapping
LLM behaviour onto human cognitive categories.

\end{description}

\subsection*{A.7\quad Conclusions}

\textbf{Stated conservatively:}

All 15 novels are available from Project Gutenberg and almost certainly
appear in LLM training data.  Yet the LLM's ability to quote them from
memory varies from 57\% to 0\%.  This gradient correlates with
critical-cultural prominence ($r = 0.72$, $n = 15$) and persists
within authors: among four Dickens novels, three Eliot novels, and two
Gissing novels, the within-author ranking tracks the novel's position
in modern critical culture.

We interpret this as a literary-critical manifestation of differential
LLM memorisation: canonical novels, whose passages are duplicated
across study guides, academic papers, and Wikipedia, are more
faithfully retained than obscure ones.  This interpretation is
consistent with \textcite{carlini2023quantifying}'s framework but
unverifiable without access to training data.

The parallel with \textcite{bayard2007comment}'s argument about human
non-reading is suggestive: both humans and LLMs navigate literary
culture through graduated approximation.  The mechanisms are entirely
different, but the observable gradient---from faithful quotation
of canonical works to creative pastiche of obscure ones---is
strikingly similar.


\newpage
\titleeng{Branje brez branja: kako bralna strategija, interpretativni okvir in besedilna utemeljenost oblikujejo ra\v{c}unalni\v{s}ko literarno razpravo}
\begin{abstracteng}
Bralec, ki roman prebere v celoti, in bralec, ki se osredoto\v{c}i le na klju\v{c}ne odlomke, zbereta razli\v{c}ne dokaze za razpravo.  Ko pa obe bralni strategiji izvedemo ra\v{c}unalni\v{s}ko in ustvarimo ve\v{c}glasne literarne razprave o 15 romanih, so rezultati skoraj enaki: dele\v{z} govora strokovnjakov je stabilen ne glede na strategijo.  To razhajanje med bralnimi strategijami in zbli\v{z}evanje razprav zagotavlja ra\v{c}unalni\v{s}ke dokaze za Fishovo tezo, da pomen konstruira bralec.
\end{abstracteng}
\keywordsslo{literarna interpretacija, teorija bralnega odziva, memorizacija jezikovnih modelov, besedilna utemeljenost, digitalna humanistika}

\end{document}
